% ==========================================================================
% 8. ESTIMATION   [full deck only]
% ==========================================================================
% PI directive 6: a large illustrative choice task in the instrument's wording plus a three-step visual of
% how values are calculated; the equations (utility, % of pay, equal-pay share) are on B-methods.
% Card: PME version, Scenario 1 layout (salary, location and schedule varied; workload, job security,
% workplace culture and commute time held at the respondent's own practice) [S5, S7]. Example pay: own pay
% x 1.10 and x 0.90, shown here as +10% / -10% [S6b]; the instrument showed a monthly S$ amount [S6].
\begin{frame}[t,label=M-estimation]{How we turn choices into values}
  \FrameTop % first body token (centring; never a frame subtitle)
  {\centering
  % Polish pass (DESIGN_PLAN s2, W2): card heads as on M-design (coloured bold text over a 1.2 pt rule, no
  % dark fill); cells deckcell, held box deckheld; each step a deckcard with a 0.07 cm Navy left bar and a
  % 12 pt bold Navy numeral in a hanging column; step text 10 pt; arrows deckarrow. The picture spans the
  % full 14.8 cm text width (card 0 to 8.1; steps 8.42 to 14.78).
  \begin{tikzpicture}[x=1cm,y=1cm,every node/.style={outer sep=0pt},
      lbl/.style={deckcell,anchor=north west,minimum width=2.9cm,text width=2.74cm,align=left},
      hdr/.style={anchor=base west,inner xsep=0.08cm,inner ysep=0pt,minimum width=2.55cm,
        text width=2.39cm,align=center,text=Employee,font=\FsBody\bfseries},
      opt/.style={deckcell,anchor=north west,minimum width=2.55cm,minimum height=1.0cm,text width=2.39cm},
      stepbox/.style={deckcard,anchor=north west,inner xsep=0pt,inner ysep=0.08cm,minimum width=6.36cm,
        text=Ink,font=\FsSmall}]
    % ---- card (x 0 to 8.1): label and scenario line [S5] six tasks
    \node[deckhead] (hc) at (0,0) {Illustrative choice task};
    \node[anchor=base east,inner sep=0pt,text=Muted,font=\FsSmall] at (8.1,0 |- hc.base) {Scenario 1 of 6}; % [S5] six tasks per respondent
    % header row (instrument: Attribute / Job Offer A / Job Offer B) on baseline y = -0.83, rule under the pair
    % as on M-design
    \node[anchor=base west,inner xsep=0.08cm,inner ysep=0pt,text=Ink,font=\FsBody\bfseries]
      at (0,-0.83) {Attribute};
    \node[hdr] at (2.95,-0.83) {Job Offer A};
    \node[hdr] at (5.55,-0.83) {Job Offer B};
    \draw[Employee,line width=1.2pt] (2.95,-0.98) -- (8.1,-0.98);
    \draw[Rule,line width=0.4pt] (0,-0.98) -- (2.9,-0.98);
    % varied rows (shaded in the instrument). Pay: [S6b] multipliers 1.10 and 0.90
    \node[lbl,minimum height=1.0cm] at (0,-1.05) {Monthly gross salary for this job};
    \node[opt] at (2.95,-1.05) {Your pay +10\%};
    \node[opt] at (5.55,-1.05) {Your pay \textminus 10\%};
    % location levels as on the card [S7]
    \node[lbl,minimum height=1.0cm] at (0,-2.11) {Work location};
    \node[opt] at (2.95,-2.11) {Fully onsite\\(0 days WFH)};
    \node[opt] at (5.55,-2.11) {1--2 days\\WFH/week};
    % schedule levels as on the card [S7]; 9am-6pm and 4-day are level labels
    \node[lbl,minimum height=1.0cm] at (0,-3.17) {Work schedule};
    \node[opt] at (2.95,-3.17) {Fixed hours\\(9am--6pm)};
    \node[opt] at (5.55,-3.17) {Compressed\\4-day week};
    % held rows (one value spanning A and B in the instrument) [S5]; the card's foot lines up with step 3's
    \node[deckheld,anchor=north west,inner xsep=0.08cm,inner ysep=0.04cm,minimum width=8.1cm,
      minimum height=0.89cm,text width=7.94cm] at (0,-4.23)
      {Workload, Job security, Workplace culture, Commute time: same for A and B (your current job)};
    % ---- three steps (x 8.42 to 14.78), stacked 0.30 cm apart
    \node[deckhead] at (8.42,0) {How values are calculated};
    % [S5] six choices per respondent
    \node[stepbox] (sA) at (8.42,-0.48)
      {\hspace*{0.22cm}\makebox[0.36cm][l]{\color{Navy}\fontsize{12}{14}\selectfont\bfseries 1}%
      \parbox[t]{5.66cm}{\RaggedRight Choose A or B, six times}\hspace*{0.12cm}};
    \node[stepbox] (sB) at ([yshift=-0.30cm]sA.south west)
      {\hspace*{0.22cm}\makebox[0.36cm][l]{\color{Navy}\fontsize{12}{14}\selectfont\bfseries 2}%
      \parbox[t]{5.66cm}{\RaggedRight Estimate how pay and each FWA shift choices}\hspace*{0.12cm}}; % [S5] step label, not a data value
    % EB s3b: % of pay = pay gap at which the model predicts a 50:50 choice (indicative, S6 and EB s9);
    % equal-pay share = invlogit of the FWA coefficient, 50% = no preference, read within each side
    \node[stepbox] (sC) at ([yshift=-0.30cm]sB.south west)
      {\hspace*{0.22cm}\makebox[0.36cm][l]{\color{Navy}\fontsize{12}{14}\selectfont\bfseries 3}%
      \parbox[t]{5.66cm}{\RaggedRight Translate the average effect into:\\[0.04cm] % [S5] step label, not a data value
      \textbf{\% of pay} (indicative): the pay gap that makes A and B equally attractive\\[0.04cm]
      \textbf{Equal-pay share}: how often the FWA option wins if pay is equal; 50\% = no preference}% [EB s3b]
      \hspace*{0.12cm}};
    \foreach \s in {sA,sB,sC} {\fill[Navy] (\s.north west) rectangle ([xshift=0.07cm]\s.south west);}
    \draw[deckarrow] ([xshift=0.33cm]sA.south west) -- ([xshift=0.33cm]sB.north west);
    \draw[deckarrow] ([xshift=0.33cm]sB.south west) -- ([xshift=0.33cm]sC.north west);
  \end{tikzpicture}\par}
  % [S6, S6b] no equal-pay pair was shown; every gap 17-33%; EB s9: magnitudes model-based
  \Interp{No equal-pay pair was shown; observed pay gaps were about 17--33\%,\\ so rankings are firmer than prices.}
  % card wording: questionnaire/FWA_Study_-_Employee_DCE.docx Section D; employer card: Candidate A / Candidate B
  \SourceNote{\textbf{Source:} PME questionnaire, Section D. Employers saw the same layout with Candidates A and B.}
  \Backups{\BL{B-methods}{Equations}\BL{B-eqpay}{Equal pay}\BL{B-validity}{Validity}}
\end{frame}
% note facts: [S5] six choices, employer version with two candidates; EB s3b: equal-pay share needs no pay
% slope, 50% = no preference, read within each side; [S6] 17 to 33 per cent; EB s9: magnitudes model-based;
% [V1-V5] stated choices predict direction and ranking well, levels less well
% implication, spoken: partners can rely on rankings; per cent of pay is a rough size, never a pay cut
\note{[Full deck only]
This is what a PME saw; employers saw two candidates. % [S5]
A statistical model estimates how pay and each arrangement shift the six choices. % [S5]
We report two numbers. Per cent of pay is the pay gap that makes the two options equally attractive. The equal-pay share is how often the flexible option wins at equal pay; 50 per cent means no preference. % [EB s3b]
We acknowledge that no pair had equal pay and gaps ran from 17 to 33 per cent, so per cent of pay extrapolates. % [S6, EB s9]
MOM, SNEF and NTUC can rely on the rankings; per cent of pay is a rough size, never a proposed pay cut. % [V1-V5]
Next, what PMEs and firms do today.}

% ==========================================================================
% 9. DESCRIPTIVE EVIDENCE   [full deck only]
% ==========================================================================
% F_landscape_m (fig_v2.py, redrawn 13.0 x 5.0 cm): PME use and take-up where offered, dagger ranges
% [D1, D2x]; employers' typical PME WFH days [D3]. Claims qualified as surveyed samples (COMMENTS slide 9).
\begin{frame}[t,label=M-descriptive]{Hybrid is common; compressed weeks are rare}
  \FrameTop
  \FigC{F_landscape_m}
  % [D3] 62% of surveyed employers call 1-2 WFH days the typical PME arrangement;
  % [D1] compressed-week use 6.5% [4.3% excl. Lucid], display 4-6% with dagger
  \Interp{Among surveyed firms, \Pend{62\%} call 1--2 WFH days typical for PMEs;\\ only 4--6\%\dg\ of surveyed PMEs use a compressed week.}
  % [D1] dagger range: all respondents and excluding one online panel; [D1, D3] PME and employer surveys
  \SourceNote{\textbf{Note:} \DgShort\ \textbf{Source:} this study's PME and employer surveys (Aug--Oct 2026).}
  \Backups{\BL{B-landscape}{Landscape}\BL{B-setup}{Pathways}}
\end{frame}
% note facts: [D1] any FWA 73% [69%], about seven in ten; [D3] 62%; [D1] compressed 4-6%; [M16] 20.8% of
% firms offer it, about one firm in five; [D5x] company-wide policy 43% [39%] vs individual request 30% [26%]
% implication, spoken: SNEF could share written FWA policy templates with member firms
\note{[Full deck only]
Flexible work is already common among the PMEs we surveyed: about seven in ten use at least one arrangement, and \Pend{62 per cent} of surveyed firms call one to two WFH days typical for their PMEs. % [D1, D3]
The compressed week is the exception: four to six per cent of PMEs use one, and \Pend{about one firm in five} offers it. % [D1, M16]
These are our samples, not national rates.
More arrangements come through company policy than individual requests, so SNEF could share written FWA policy templates with member firms. % [D5x]
Is the compressed week rare because PMEs place little value on it, or because employers hesitate? PMEs first.}

% ##########################################################################
% Main slides from redesign/v3/W3_frames.tex
% ##########################################################################
\begin{frame}[t,label=M-pmes]{PMEs strongly prefer hybrid work}
  % F_emp_wtp, true size 13.0 x 5.0 cm: % of monthly pay a PME would give up for each level, 95% CIs.
  % Rows: fully remote vs own commute 12.7 [9.8, 15.6] [M1]; 3-4 WFH days 9.5 [E2]; 1-2 WFH days 7.4 [E1];
  % compressed week 7.4 [E5]; staggered hours -0.7 [E4]; job sharing 6.0 [E6] and part-time 3.9 [E7],
  % hours change too; shaded band: WFH option abroad 4-8% [E16]. Fixer pass 7 Oct 2026 (fact_audit s3 item 7):
  % the figure prints these status-I values as whole numbers, 13, 10, 7, 7, -1, 6, 4 (EB: quote rounded, as about)
  \FrameTop
  \FigC{F_emp_wtp}
  % [E1, E2] 1-2 days 7.4, 3-4 days 9.5 -> about 7-10%; [S6] pay gaps 17-33% by design (mandated wording kept)
  \Interp{The model implies about 7--10\% of pay for hybrid work:\\
    indicative, not a proposed pay cut; observed pay gaps were about 17--33\%.}
  % Note: F_emp_wtp's former footnote, verbatim, for the starred job-sharing and part-time rows [E6, E7].
  % Polish pass: "Conditional logit" leaves the source line (DESIGN_PLAN s2 W3; M-estimation and B-methods
  % carry the method). One line at 7.5 pt, 406 of 421 pt: the source names the survey, its size and month,
  % as on M-data; the exact field dates 4-21 Aug 2026 stay on F_sample_m (M-data) and its speaker note. Fixer pass
  % 7 Oct 2026 (employer-01, mom_ps-06, worker-04): never "SMU--SNEF survey" until the PI confirms roles (G0).
  % [S1] 1,020 PMEs, fielded 4-21 Aug 2026, online panels; 95% CIs clustered by respondent (EB s2 header)
  \SourceNote{\textbf{Note:} * Hours change too: not a price of flexibility. Lines: 95\% CIs. \textbf{Source:} PME survey (n = 1,020), Aug 2026 (online panels).}
  \Backups{\BL{B-emp-all}{All estimates}\BL{B-remote-commute}{By commute}}
\end{frame}
% note facts: [E1, E2] 7 to 10; [E16] at or above the upper end of estimates abroad; [E5] about 7;
% [E4] about zero; [M1] about 13 vs own commute; [S6] 17 to 33; [E17] bracketed Q&A cue only, illustrative
\note{[Short path] First, what surveyed PMEs value: the share of monthly pay they would give up for each arrangement, against fully onsite work.
Hybrid work, one to two or three to four WFH days, is worth about 7 to 10 per cent of pay, at or above the upper end of estimates abroad. A compressed week, about 7; staggered hours, about zero; fully remote, against one's own commute, about 13. % [E1, E2, E16, E5, E4, M1]
We acknowledge that the pay gaps shown were about 17 to 33 per cent, so these levels are indicative, not a proposed pay cut. % [S6]
For firms and SNEF, hybrid options can help attract PMEs alongside pay, not instead of it.
Yet one group chooses WFH less often.
[If asked for dollars: at a median S\$6,000 a month, about S\$400--600; illustrative only.]} % [E17] bracketed cue, not timed

% ==========================================================================
% 11. MECHANISMS: CAREER CONCERN AND WFH CHOICES (association)   [short path]
% ==========================================================================
\begin{frame}[t,label=M-why-pmes]{Career concerns track lower demand for WFH}
  % F_why_pmes, true size 13.0 x 5.0 cm: flexible job's % of choices at equal pay, 95% CIs, by career belief
  % (QID51_1 agree 4-5, n = 455, vs others, n = 565) [M4]; legend "No expected penalty / Expected penalty".
  % 1-2 WFH days 52 [47, 57] vs 64 [60, 69]; 3-4 WFH days 52 vs 69; remote vs a 15-45 min commute 55 vs 72;
  % location test p < .001 [M4]; compressed week 59 [54, 64] vs 58 [52, 63], p = .70 [M5]; 50% line (EB s3b)
  \FrameTop
  \FigC{F_why_pmes}
  % [M6] items 3-4: belief asked after the choice tasks; gap smaller excl. one panel, WTP p = .048 there
  \Interp{Consistent with a visibility concern, but asked after the choices\\
    and weaker outside one panel.}
  % Note: F_why_pmes's former footnote. [M4] penalty = QID51_1 agree 4-5; "Associations" kept. Its second
  % clause, gap smaller excl. one panel [M6], is not repeated: the Interp above states it, weaker outside
  % one panel. Equal-pay shares: the figure header names the measure. One line, 412 of 421 pt.
  % [S1] 1,020 PMEs, 4-21 Aug 2026; [M4] fully interacted model, equal-pay shares, 95% CIs
  \SourceNote{\textbf{Note:} Penalty: agrees flexible hours would hurt their career. Associations; lines: 95\% CIs. \textbf{Source:} PME survey (n = 1,020), Aug 2026.}
  \Backups{\BL{B-why-pmes}{All splits}\BL{B-career-who}{Career risk}}
\end{frame}
% note facts: [M4] 52 vs 64, 1% level; [M5] compressed week 59 vs 58, p = .70, no detectable difference;
% [M6] post-choice item, about 40% of the gap could be noisier choices, weaker excl. one panel
\note{[Short path] Surveyed PMEs who expect shorter or more flexible hours to hurt their career choose the job with one to two WFH days 52 per cent of the time at equal pay, against 64 per cent for others, statistically significant at the 1 per cent level. For the compressed week, we detect no difference. % [M4, M5]
This is consistent with a visibility concern. We acknowledge that the career question was asked after the choices, that part may be noisier answers, and that it is weaker outside one panel. % [M6]
For firms and NTUC, appraisal by output could be part of any hybrid pilot.
What do employers weigh in hiring?}

% ==========================================================================
% 12. RESULT 2: WHAT EMPLOYERS RANK, IN HIRING   [short path]
% ==========================================================================
\begin{frame}[t,label=M-employers]{Employers rank fully remote last in hiring}
  % F_org_eqpay, true size 13.0 x 5.0 cm: share of hiring choices the flexible candidate wins at equal salary,
  % 95% CIs; no % of salary values (mandate; EB s8). Appended sample (n = 397), appended_clogit.xlsx main,
  % invlogit(b) with delta-method CIs. Rows: compressed week 47 [41, 53] [Q5];
  % staggered hours 45 [40, 50] [Q4]; 1-2 WFH days 40 [35, 45] [Q1]; 3-4 WFH days 31 [26, 36] [Q2];
  % fully remote 23 [18, 28] [Q3]; grey benchmark: >20% turnover-risk candidate 33.6 [28, 39] [Q8];
  % F_org_eqpay redrawn for the appended sample (fig_r5.py, 14 Oct 2026)
  % 50% = no preference (EB s3b); ranking [R10]
  \FrameTop
  \FigC{F_org_eqpay}
  % EB s3b: 50% = indifferent, equal-pay share uses no salary coefficient; [R13] pay is a weak lever for employers
  % line 2, fixer pass 7 Oct 2026 (employer-02): MANDATE information/learning problem; [S5; EB s9] hiring choice
  % between candidates, not a request and not an operating-cost measure; the DCE cannot separate these; EB s3b 50%
  \Interp{Ranking, not a salary price: shares of hiring choices at equal pay; 50\% = no preference.\\
    Reservations may reflect operating needs, uncertainty or a candidate signal.}
  % Note: F_org_eqpay's former footnote, verbatim. [S6] no shown pair had equal pay, gaps 17-33%; EB s3b model-based share.
  % Source: the respondent wording, people involved in PME hiring or manpower decisions, stays on M-design
  % and M-data; here "employer survey, n = 397, Aug-Oct 2026" keeps the block on one line, 408 of 421 pt.
  % [S2] 397 respondents (appended sample) involved in PME hiring/HR/manpower decisions; [S5] hiring choice
  \SourceNote{\textbf{Note:} Lines: 95\% CIs. Model-based shares: no pair in the survey had equal pay. \textbf{Source:} employer survey (n = 397), Aug--Oct 2026.}
  \Backups{\BL{B-org-yardstick}{Yardstick}\BL{B-org-salary}{Salary lever}}
\end{frame}
% note facts: [S2, S5] hiring choice between equally qualified candidates; EB s3b 50 = no preference;
% [Q5] 47; [Q1] 40; [Q2] 31; [Q3, R10] 23, ranked last (appended, n = 397); [Q16] adjacent tests NOT in the
% appended_* outputs (303-firm values: compressed vs 1-2 days p = .22, 1-2 vs 3-4 days p = .061, 3-4 days vs remote p = .006);
% [R13] employers picked the cheaper candidate only 57% of the time, a weak salary response
% [Q8] benchmark: candidate with a >20% turnover risk wins 33.6% (28-39, appended), the grey row of F_org_eqpay (bracketed cue, fixer pass)
% [D3] 62% of surveyed employers call 1-2 WFH days the typical PME arrangement (employer-02; MANDATE information/learning)
\note{[Short path] Now employers. These are hiring choices between two equally qualified candidates, not responses to a staff request. Each bar is how often the flexible candidate wins at equal salary; 50 per cent means no preference. % [S2, S5]; EB s3b
A compressed-week candidate wins 47 per cent: no detectable hiring penalty. One to two WFH days win 40, three to four days 31, and fully remote 23. \Pend{Only the drop to fully remote is statistically significant, at the 1 per cent level.} % [Q5, Q1, Q2, Q3, R10] appended; [Q16] adjacent tests not re-run
[For scale: the grey row, a candidate with a high turnover risk, wins 33.6 per cent.] % [Q8] benchmark on F_org_eqpay (appended)
\Pend{Salary moved choices only weakly}: a ranking, not a price. % [R13] model-free share not in appended_* outputs
Most surveyed firms already use hybrid work: \Pend{62 per cent} call one to two WFH days typical for PMEs. So a hiring reservation may reflect role needs, uncertainty or a candidate signal, not simply opposition to hybrid work. % [D3]
SNEF and member firms could start pilots near the top of this ranking.
Why are some more reserved?}

% ==========================================================================
% 13. MECHANISMS: WHERE EMPLOYER RESERVATIONS SIT (association)   [short path]
% ==========================================================================
\begin{frame}[t,label=M-where]{Employer reservations vary with role needs and concerns}
  % Title: 13.9 cm at 14.4 pt bold (layout_sim Helvetica metrics) < 14.8 cm, one line, so the mandated title is kept.
  % F_where_r5, true size 13.0 x 5.0 cm: two contrasts, % of hiring choices at equal pay, 50% line (EB s3b).
  % 1-2 WFH days: most PME roles need presence (n = 98) 29% vs only some or few roles (n = 205) 49% [M10, Q10];
  % fully remote: creativity concern (n = 162) 19% vs not worried (n = 141) 39% [M10, Q11];
  % location tests p = .006 and .003, common salary slope [M10]
  % Appended sample (appended_clogit.xlsx physical_1/2, creative_concern_1/0; invlogit, delta-method CIs):
  % 1-2 WFH days: most roles need presence (n = 137) 27% [19, 34] vs only some or few (n = 260) 48% [42, 55];
  % fully remote: creativity concern (n = 201) 18% [11, 25] vs not worried (n = 196) 29% [21, 37].
  % F_where_r5 redrawn from these values (MOM_presentation/scripts/make_figs_appended.py). TODO(appended): the
  % pooled location tests (p) and the 20-split q are not in the appended_* outputs
  \FrameTop
  \FigC{F_where_r5}
  % [M12, Q15] BH q = .058 across 20 splits; EB s9: both items asked after the choices (Block 5); EB s9: no
  % group-specific salary-slope check yet for these splits, so "preliminary association under a common pay-slope model"
  \Interp{Preliminary associations under a common pay-slope model;\\
    both items asked after the choices; \Pend{q \ap .06 across 20 splits}.}
  % Note: F_where_r5's former footnote, verbatim, the filled and hollow marker key [M10]. One line, 408 of 421 pt.
  % CIs: separate per-group clogits, EB M10 drawn CIs (fixer pass 7 Oct 2026, editor-04)
  % [S2] 397 respondents (appended sample); [M10] filled = firms with the need or concern
  \SourceNote{\textbf{Note:} Filled: firms with the need or concern; hollow: other firms. Lines: 95\% CIs. \textbf{Source:} employer survey (n = 397), Aug--Oct 2026.}
  % employer-07: employers' own view (B-org-view returns to M-employers and M-where)
  \Backups{\BL{B-why-org}{Tests}\BL{B-org-concerns}{Concerns}\BL{B-org-view}{Employer view}}
\end{frame}
% note facts: [M10, Q10] 29 vs 49; [M10, Q11] 19 vs 39; [M12, Q15] p <= .006 individually, q = .058
% across 20 splits, i.e. the 10% level after correction; EB s9 post-choice items and the pending
% group-specific salary-slope check; [P11] Job Redesign+ grant; [B2] 67% say FWAs help attract and retain talent;
% [B1] biggest challenges coordination 25%, fairness to staff who cannot have FWAs 23%, productivity 22% (employer-07)
\note{[Short path] Where most PME roles need on-site presence, a candidate wanting one to two WFH days wins 27 per cent at equal pay, against 48 elsewhere. Where firms worry that creativity suffers apart, a fully remote candidate wins 18 per cent, against 29. % [M10, Q10, Q11] appended sample
\Pend{Each is statistically significant at the 1 per cent level, but only at the 10 per cent level after correcting for 20 splits.} Both items were asked after the choices, and a check letting pay sensitivity differ across firms is pending: these are preliminary associations. % [M12, Q15]; EB s9
\Pend{Two in three} surveyed employers say FWAs help attract and retain talent; their biggest hurdles are coordination, fairness to staff without FWAs, and productivity. % [B2, B1] not in appended_* outputs
Firms, with Job Redesign+ support, could first review which tasks need presence. % [P11]
Next, each side against its own 50 per cent.}

% ##########################################################################
% Main slides from redesign/v3/W4_frames.tex
% ##########################################################################
\begin{frame}[t,label=M-meet]{Each side, against its own 50\%} % [s3b] 50% = no preference, read within each side only (Lead decision 3)
  % F_meet_r5 (fig_r5.py, 14.0 x 4.7 cm; integrator pass 7 Oct 2026: the header on one
  % line gives the slide 0.30 cm of slack): two visibly separate panels, each with its
  % own axis and its own 50% reference line [s3b]; same row order in both panels;
  % fixer pass 7 Oct 2026 (mom_ps-09): one x-range per panel, PMEs 40-80 and employers 10-70, so the panels share
  % no ruler; footnote (editor-01): zones A-C are descriptive labels [Q16]; PMEs' remote vs own commute [M1];
  % header printed in the figure: "Compare each side with its own 50%; do not
  % compare distances across sides" (mandate, M-meet). Zone labels A/B/C inside are
  % descriptive labels [Q16]. Model-based equal-pay shares with 95% CIs:
  %   left, PMEs:       compressed week 58 (55-62) [Q5]; 1-2 WFH days 58 (55-62) [Q1];
  %                     3-4 WFH days 61 (57-64) [Q2]; fully remote 64 (61-68, vs own commute) [M1, Q3]
  %   right, employers (appended, n = 397): compressed week 47 (41-53) [Q5]; 1-2 WFH days 40 (35-45) [Q1];
  %                     3-4 WFH days 31 (26-36) [Q2]; fully remote 23 (18-28) [Q3]
  %   F_meet_r5 redrawn for the appended sample (fig_r5.py, 14 Oct 2026)
  % "1--2", "3--4 WFH days" are attribute levels [S7], not data values.
  \FrameTop
  \FigC{F_meet_r5}
  % compressed-week wording: mandate "Evidence and language rules"; [Q5] 47% (41-53, appended) includes 50
  \Interp{Compressed week: most promising common-ground test;\\ evidence moderate; operating costs unmeasured.}
  % Note = F_meet_r5's former footnote, verbatim (polish pass 7 Oct 2026): [s3b] equal-pay shares, 95% CIs;
  % [Q16] zones A-C are descriptive labels; [M1] PMEs' fully remote is measured against each PME's own commute.
  % [S1] 1,020 PMEs, 4-21 Aug 2026; [S2] 397 employer respondents (appended sample)
  \SourceNote{\textbf{Note:} \% of choices at equal pay; 95\% CIs. Zones A--C: descriptive labels. PMEs' remote vs own commute.
    \textbf{Source:} choice experiments with 1,020 surveyed PMEs and 397 employer respondents (Aug--Oct 2026).}
  \Backups{\BL{B-eqpay}{Equal pay} \BL{B-language}{Wording}}
\end{frame}
\note{[Short path] Two panels, two scales. Read each against its own 50 per cent line, meaning no preference; do not compare distances across sides. % [s3b; Lead decision 3; mandate M-meet]
The compressed week, the same hours over four longer days, is the most promising common-ground test: PMEs pick it 58 per cent of the time at equal pay, and employers show no detectable hiring penalty. Evidence is moderate; operating costs are unmeasured. % [S7, Q5] employers 47% (41-53, appended)
One to two WFH days depend on the role. Employers pick candidates wanting three to four days or fully remote well under half the time; these need testing before scaling. % [Q10, Q2, Q3] 31% and 23% (appended)
SNEF and volunteer firms could test the compressed week first. So why is it rare? % implication (who could act) and transition to M-puzzle (short path)
[If asked: the zones are descriptive labels; \Pend{only the step from three to four days to fully remote is statistically significant at the 1 per cent level}.] % [Q16] p = .006, 303-firm value; not re-run
[Full deck: go to M-groups first and ask: does demand differ across PMEs?]}

% ==========================================================================
% 15. HETEROGENEITY ACROSS PME GROUPS   [full deck only]
% ==========================================================================
\begin{frame}[t,label=M-groups]{No subgroup difference detected after correction}
  % F_groups_m (fig_v2.py, 13.0 x 5.0 cm; compressed-week column removed, ROUND5_SPEC s4):
  % PME equal-pay shares with 95% CIs by gender, age (19-39 / 40+) and caregiving (none / any),
  % for 1-2 WFH days, 3-4 WFH days and fully remote (vs a long commute); 50% reference [s3b].
  %   [Q14b] every plotted group's CI lies above 50% for the three WFH levels (lowest bound: men, 1-2 days, 51.2)
  %   [Q14] women 63/67/73% vs men 56/56/63% (location p = .047, nominal only); age p = .30; caregiving p = .53
  %   [E18] 198 pairwise contrasts, minimum BH q = .70; WTP joint tests p = .43 / .09 / .66
  \FrameTop
  \FigC{F_groups_m}
  % [E18, M7] null results are not equivalence (mandate: "Non-rejection is not equivalence")
  \Interp{Non-rejection is not equivalence; groups may still differ.}
  % Note (polish pass 7 Oct 2026): the 95% CI convention (formerly in the figure's axis caption and this
  % source line, [s3b]); "Fully remote: vs a >45 min commute", moved out of the figure's column title
  % ([E3] main model: fully remote vs onsite with a long commute; [S8] long = >45 min); and the figure's
  % former "* Nominal" footnote, verbatim ([Q14] the starred p = .047 is nominal; [E18] 198 contrasts, BH).
  % [S1] 1,020 PMEs; [E18] BH correction over 198 contrasts; [M7] 7 primary splits
  \SourceNote{\textbf{Note:} Lines: 95\% CIs. Fully remote: vs a \textgreater 45 min commute. * Nominal: none of 198 subgroup contrasts survives multiple-testing correction.
    \textbf{Source:} PME choice experiment (1,020 surveyed PMEs).}
  \Backups{\BL{B-emp-groups}{All groups} \BL{B-why-pmes}{All splits}}
\end{frame}
\note{[Full deck only] In every group shown, the interval for choosing a WFH job at equal pay lies above one half. % [Q14b]
None of 198 group contrasts survives correction for multiple testing, and we detect no differences by income, organisation type or tenure. % [E18, M7]
We acknowledge that no difference detected is not proof that groups are alike; mothers of young children are a small sample. % [M7] n = 108
Only the career-concern split survives, as we saw earlier. % [M7, M4] a belief, not a demographic group
MOM, SNEF and NTUC could therefore treat PME demand for WFH as broad, not niche. Why, then, is the compressed week rare? % implication and transition to M-puzzle
}

% ==========================================================================
% 16. THE COMPRESSED-WEEK PUZZLE   [short path]
% ==========================================================================
\begin{frame}[t,label=M-puzzle]{The compressed week: valued, rarely offered}
  % F_puzzle_r5 (fig_r5.py, 13.0 x 5.0 cm), simplified to the puzzle only:
  %   PMEs choose the compressed-week job 58% of the time at equal pay (55-62) [Q5]
  %   firms offering it: about 1 in 5 (63 of 303 = 20.8% [16.6, 25.7]) [M16] TODO(appended): not in appended_* outputs
  %   offered by 26.5% of firms with a written FWA policy vs 9.7% without (association) [M17]
  %   take-up where available: about 1 in 3 (65 of 210 = 31.0%) [M22]
  % "compressed week" = the same total hours over four longer days [S7]
  \FrameTop
  \FigC{F_puzzle_r5}
  % pilot framing: COMMENTS slide 16; mandate (voluntary, evaluated; not costless)
  \Interp{This supports a voluntary, evaluated pilot, not a claim that adoption is costless.}
  % Note (polish pass 7 Oct 2026): the figure's former footnote ("+ All respondents / excluding one online
  % panel. Surveyed samples, unweighted.") with the deck's one dagger wording \DgShort [S1, S2];
  % [M20] cheaper option in 83.2% of PME pairs (1,502/1,805) and 90.8% of employer pairs (327/360)
  \SourceNote{\textbf{Note:} \DgShort\ Surveyed samples, unweighted.
    Design: where only one option had it, the compressed week was the cheaper one in 83\% of PME and \Pend{91\%} of employer pairs.}
  \Backups{\BL{B-puzzle}{Offering} \BL{B-pilot}{Pilot checklist}}
\end{frame}
\note{[Short path] PMEs value the compressed week and employers show no detectable hiring penalty, yet only \Pend{about one surveyed firm in five} offers it. % [Q5, M16] 58% in the figure
Firms with a written FWA policy offer it more often, \Pend{27 against 10 per cent}, though a policy may follow offering rather than lead it. % [M17] association
Where it is available, about one surveyed PME in three takes it up. % [M22]
We acknowledge that the design usually made it the cheaper option, and a hiring experiment cannot measure rostering, fatigue or service coverage. % [M20]
So the case is for a voluntary, evaluated pilot, not a roll-out: SNEF could recruit volunteer firms to put it on their written menu, with NTUC co-designing safeguards. What happens when PMEs ask? % implication (who could act) and transition to M-barriers
}

% ==========================================================================
% 17. BARRIERS: CAREER CONCERN, AWARENESS, REQUEST OUTCOMES   [short path]
% ==========================================================================
% Title: ROUND5_SPEC s3 row 17 wording, restored by the integrator (7 Oct 2026). W4 had shortened it under the
% superseded 48-character lint rule; title_fit.py measures it at 14.37 cm bold (13.36 regular) of 15.4 cm: one line.
\begin{frame}[t,label=M-barriers]{We see outcomes for requesters, not why others do not ask}
  % F_barriers_r5 (fig_r5.py, 13.0 x 5.0 cm): three facts only, each with its own base; left header
  % Not conditional on asking, right header Reported requesters (fixer pass 7 Oct 2026: the left header was
  % All surveyed PMEs, the wrong base for card b; fact_audit WRONG 3, mom_ps-05, worker-10):
  %   (a) all surveyed PMEs: 37-45%+ expect shorter or more flexible hours to hurt their
  %       career progression (44.6% all / 37.1% excl. Lucid) [M23, B5]
  %   (b) junior PMEs (employed): only 25-29%+ familiar with the Guidelines (29% all / 25% excl. Lucid) [D6];
  %       caption adds all levels, employed PMEs N = 1,004: 44-55%+ (55% all / 44% excl. Lucid) [D6] (mom_ps-05;
  %       a ROUND5_SPEC row 17 adaptation logged in ADJUDICATION.md)
  %   (c) among PMEs who report a formal request: 4-7%+ report an outright refusal
  %       (4.4% all / 7.2% excl. Lucid), a selected lower bound [D7, D7b; item misread D9]
  % + = dagger: range across all respondents and excluding one online panel.
  \FrameTop
  \FigC{F_barriers_r5}
  % bases separated (mandate M-barriers; COMMENTS slide 17): [M23] all PMEs; [D7] reported requesters only; [B6] no why-not item
  \Interp{Career concern covers all surveyed PMEs; refusals cover only reported requesters.\\ We cannot say why others did not ask.}
  % Note (polish pass 7 Oct 2026): the figure's former footnote ("+ All respondents / excluding one online
  % panel (grey: the gap between the two)"), with the dagger stated once as \DgShort [B5, D6, D7b];
  % [S1] PME survey 4-21 Aug 2026; [D9] request item partly misread ("a selected lower bound" in the figure; the note)
  \SourceNote{\textbf{Note:} \DgShort\ Grey: the gap between the two.
    \textbf{Source:} PME survey (Aug 2026).}
  \Backups{\BL{B-career-who}{Who expects} \BL{B-approvals-tab}{Approvals}}
\end{frame}
\note{[Short path] Three facts, on different bases. First, across all surveyed PMEs, 37 to 45 per cent expect shorter or more flexible hours to hurt their careers. % [M23, B5]
After accounting for who answered, parents of young children are somewhat more likely to expect a career penalty. % [M26] EB wording for NTUC, no number
Second, only a quarter to three in ten surveyed junior PMEs are familiar with the Guidelines. % [D6]
Third, among PMEs who report a formal request, four to seven per cent report an outright refusal, or about seven to ten per cent once arrangements that apply automatically are set aside, because the request item was partly misread: a selected group, so both are lower bounds. % [D7, D7b, D9]
We cannot see why PMEs who need an arrangement do not ask. MOM's surveys could add that question, and any pilot needs career safeguards that NTUC could help design. % [B6] implication (who could act)
One guardrail before the options. These are stated choices from unweighted online panels; rankings are firmer than prices; and the employer splits by role and concern remain preliminary associations. % 20-second inferential guardrail (COMMENTS slide 21; mandate M-limits); transition to M-options
[If asked: familiarity goes with more expected career risk, not less.] % [M25]
[If asked: among those who asked, 93 to 96 per cent were approved fully or partly, broadly in line with MOM's figure, though definitions differ.] % [D7c, D7, P3]
[If asked: MOM finds about 6 per cent of FWA users report a setback; that is experience, ours is expectation.] % [P15, M23]
[Full deck: M-robust and M-limits come next; the guardrail can be shortened there.]}

% ##########################################################################
% Main slides from redesign/v3/W5_frames.tex
% ##########################################################################
\begin{frame}[t,label=M-robust]{Rankings are more credible than levels}
  % T_robust_main (make_tables.py; do not edit): check | concern | result, 6 rows.
  % Row fact IDs are in the table file: [E15, E2] excl. one panel <=2 points;
  % [E15, RB2] outliers, trim, same-side <=1 point; [RB2] position term <=2;
  % [R4, EB s9] version clusters, 60 pairs, staggered p = .15; [RB1, E11, R10]
  % rank pay model fits about as well, same ranking.
  % Printed in the shared main-table style, \MainTableStyle = 10 pt (ROUND5_SPEC s3: tables >= 9 pt;
  % polish pass G8, 7 Oct 2026; was 9 pt); \Tab would force the 8.5 pt backup \TableFont.
  \FrameTop % first body token (centring; never a frame subtitle)
  {\centering\MainTableStyle\input{tables/T_robust_main}\par}
  % [S6, S6b] pay gaps 17-33% by design (displayed about 16-34%; fixer pass 7 Oct 2026, editor-13): % of pay values are model-based;
  % [EB s9] 10 versions x 6 tasks = 60 distinct choice situations per survey;
  % [V1-V5] no cited validation study tests employer hiring choices (B-validity)
  \Interp{Pay gaps were 17--33\% by design, so levels in \% of pay are model-based;\\ 60 choice situations per side; employer hiring-task validity untested.}
  % [S1] 1,020 PMEs, 4-21 Aug 2026; [S2] 397 employer respondents (appended sample)
  \SourceNote{\textbf{Source:} choice experiments with 1,020 surveyed PMEs and 397 employer respondents (Aug--Oct 2026).}
  \Backups{\BL{B-validity}{Validity} \BL{B-emp-robust}{Checks}}
\end{frame}
\note{[Full deck only] Which results survive changes to the analysis? % [E15, RB2]
Dropping the weaker panel, implausible salaries or inattentive respondents, or adding a position term, moves PME values by two points of pay or less; the rankings hold. % [E15, RB2, E11, R10]
But the pay gaps shown were about 17 to 33 per cent, so values in per cent of pay are model-based. % [S6, S6b]
We acknowledge two further limits: each side saw only 60 choice situations, so intervals are likely too narrow, and no cited study validates employer hiring choices like these. % [EB s9; V1-V5]
MOM, SNEF and NTUC can therefore act on rankings; prices need a wave with small pay gaps. How do the zones compare? % implication (who could act) and transition to M-discussion
}

% ==========================================================================
% 19. DISCUSSION: A SYNTHESIS MAP BY ZONE   [full deck only]
% ==========================================================================
\begin{frame}[t,label=M-discussion]{What the evidence says about each zone}
  % One centred three-row map (ROUND5_SPEC s3 row 19; COMMENTS slide 19: not a
  % scatter). Rows = zones (A Common ground, B Depends on the role, C Far apart;
  % descriptive labels, Q16); columns = PME side, employer side (hiring), what can
  % block use, policy move. No numbers are printed; "1--2" and "3--4" are
  % attribute levels [S7].
  %   A: PMEs value it [Q5 58%]; no detectable hiring penalty [Q5 47% appended, M19];
  %      rarely offered [M16 20.8%]; operating costs unmeasured (mandate) -> Option A
  %   B: PMEs value it [Q1 58%], less so with an expected career penalty [M4, M6];
  %      employer reservations where most roles need presence [M10, Q10], a
  %      preliminary association under a common pay-slope model [M12; EB s9 l.183]
  %      -> Option B, incl. flexi-load [P11]
  %   C: PMEs value them, remote most [Q2, M1, E11]; strong reservations, remote
  %      ranked last [Q2, Q3, R10]; creativity concern [M10, M11] (preliminary,
  %      as B) -> Option C
  %   Strip: no administrative request counts [P2]; no why-not item [B6] -> Option D. Fixer pass 7 Oct 2026 (mom_ps-02):
  %   not in published data, not unmeasured, until MRSD confirms whether its survey already asks it (PI gate)
  % Polish pass, 7 Oct 2026 (DESIGN_PLAN s2 W5; G7): the M-framework grammar. One full-row Soft band per
  % zone, no per-cell boxes and no navy bars; a deckdisc letter, the zone name in bold Ink and the
  % arrangement on the next line; heads as on M-framework; every cell text in Ink at \FsSmall = 10 pt;
  % shared deck styles only (deckcolhead, deckbody, deckdisc). Wording unchanged; line breaks only.
  % Columns (cm): disc 0.20-0.75 | zone 0.90-3.75 | PMEs 3.95-6.05 | employers 6.25-8.95 |
  % barriers 9.15-11.65 | policy move 11.85-14.65. Bands 1.45 cm tall, 0.15 cm apart; cells centred
  % vertically in their band; every cell holds at most 3 lines.
  \FrameTop % first body token (centring; never a frame subtitle)
  {\centering
  \begin{tikzpicture}[x=1cm,y=1cm,
      zc/.style={deckbody,anchor=west,font=\FsSmall}]
    % column heads, as on M-framework
    \node[deckcolhead,text=Ink] at (0.20,0.15) {Zone};
    \node[deckcolhead,text=Employee] at (3.95,0.15) {PMEs};
    \node[deckcolhead,text=Employer] at (6.25,0.15) {Employers};
    \node[deckcolhead,text=Ink] at (9.15,0.15) {Barriers};
    \node[deckcolhead,text=Ink] at (11.85,0.15) {Policy move};
    % one full-row band per zone
    \fill[deckcard] (0,0) rectangle (14.8,-1.45);
    % final fixer (7 Oct 2026): the note now prints at 8.0 pt on 2 lines; bands 0.10 cm apart (were 0.15) and
    % the dashed row 0.25 cm below them (was 0.30) give back 0.15 cm, so the table no longer touches the title bar
    \fill[deckcard] (0,-1.55) rectangle (14.8,-3.00);
    \fill[deckcard] (0,-3.10) rectangle (14.8,-4.55);
    % row A: common ground
    \node[deckdisc] at (0.475,-0.725) {A};
    \node[zc,text width=2.85cm] at (0.90,-0.725) {\textbf{Common ground}\\ Compressed week};
    \node[zc,text width=2.10cm] at (3.95,-0.725) {Value it};
    \node[zc,text width=2.70cm] at (6.25,-0.725) {No detectable hiring penalty};
    \node[zc,text width=2.50cm] at (9.15,-0.725) {Rarely offered; operating costs unmeasured};
    \node[zc,text width=2.80cm] at (11.85,-0.725) {\textbf{Voluntary, evaluated pilot (A)}};
    % row B: depends on the role
    \node[deckdisc] at (0.475,-2.275) {B};
    \node[zc,text width=2.85cm] at (0.90,-2.275) {\textbf{Depends on}\\ \textbf{the role}\\ 1--2 WFH days}; % [S7] attribute level
    \node[zc,text width=2.10cm] at (3.95,-2.275) {Value it; less with career concern};
    \node[zc,text width=2.70cm] at (6.25,-2.275) {Reservations where most roles need presence\dg};
    \node[zc,text width=2.50cm] at (9.15,-2.275) {Presence needs; career concern};
    \node[zc,text width=2.80cm] at (11.85,-2.275) {\textbf{Review tasks by role; flexi-load (B)}};
    % row C: far apart
    \node[deckdisc] at (0.475,-3.825) {C};
    \node[zc,text width=2.85cm] at (0.90,-3.825) {\textbf{Far apart}\\ 3--4 WFH days;\\ fully remote}; % [S7] attribute level
    \node[zc,text width=2.10cm] at (3.95,-3.825) {Value them; remote most};
    \node[zc,text width=2.70cm] at (6.25,-3.825) {Remote ranked last; strong reservations};
    \node[zc,text width=2.50cm] at (9.15,-3.825) {Creativity concern\dg; career concern};
    \node[zc,text width=2.80cm] at (11.85,-3.825) {\textbf{3--4-day WFH pilot; test before scaling (C)}}; % [S7] attribute level
    % strip: measurement, every zone (dashed Muted, as the M-framework barrier strip)
    \draw[Muted,dashed,line width=0.6pt] (0.02,-4.80) rectangle (14.78,-5.45);
    \node[anchor=center,inner sep=0pt,font=\FsBody,text=Ink] at (7.40,-5.125)
      {\textbf{Every zone:} no request counts, and why PMEs do not ask is not in published data \textcolor{Navy}{\textbf{(D)}}};
  \end{tikzpicture}\par}
  % dagger: [M10, M11, M12] associations from splits that use one salary slope; EB s9 l.183 free-slope check pending;
  % employer side = hiring choices between candidates [S2, S5]: carried by the A cell (hiring penalty) and the note;
  % [S1, S2] fieldwork 4-21 Aug 2026 (PMEs) and 12 Aug-2 Oct 2026 (employers)
  \SourceNote{\textbf{Note:} \dg\ Preliminary association under a common pay-slope model. \textbf{Source:} PME and employer surveys (Aug--Oct 2026).}
  \Backups{\BL{B-intl}{Trials abroad} \BL{B-puzzle}{Offering}}
\end{frame}
\note{[Full deck only] This map brings the results together and links each zone to an option. % [S7]
For the compressed week, PMEs value it and we detect no hiring penalty, yet few surveyed firms offer it, and operating costs are unmeasured; the move is a voluntary, evaluated pilot. % [Q5, M19, M16]
For one to two WFH days, reservations are concentrated where most roles need presence, a preliminary association, so firms could review tasks by role. % [M10, Q10, M12]
For three to four days and fully remote work, the two sides are far apart, so pilots come before scaling. % [Q2, Q3, R10]
In every zone, request data are missing. Which options could the partners take up? % [P2, B6]; transition to M-options
}

% ==========================================================================
% 20. OPTIONS FOR THE PARTNERS' DECISION   [short path]
% ==========================================================================
\begin{frame}[t,label=M-options]{Options for the partners' decision}
  % Decision banner (ROUND5_SPEC s3 row 20; mandate: one measurement step and one bounded practice pilot, a lead for each)
  % MANDATE: one bounded practice pilot; voluntary as on M-puzzle. Fixer pass 7 Oct 2026: mom_ps-01 (a) proposed
  % default visible, D and A, and bounded; employer-04 voluntary. Wording change vs ROUND5_SPEC row 20 logged in ADJUDICATION.md
  \FrameTop
  \DecisionBanner{Decision today: one measurement step (we propose D) and one bounded, voluntary pilot (we propose A), each with a lead.}
  \vspace{0.12cm}% polish pass, 7 Oct 2026 (DESIGN_PLAN s2 W5): was 0.04 cm
  % Three columns (COMMENTS slide 20): option | owner and first step | evidence strength. 9.5 pt.
  % Option letters link to B-optA..B-optD (detail for Q&A). No "easy", no "low cost";
  % operating cost was not measured (mandate).
  % Polish pass, 7 Oct 2026 (G8): the shared tight table size \FsTableTight = 9.5 pt, arraystretch 1.1 and
  % 2 pt of space between the option rows, so each option reads as one unit (was 1.0 and no row space).
  {\FsTableTight
  \setlength{\tabcolsep}{4pt}
  \renewcommand{\arraystretch}{1.1}
  \begin{tabularx}{\linewidth}{@{}>{\raggedright\arraybackslash}X>{\raggedright\arraybackslash}p{5.0cm}>{\raggedright\arraybackslash}p{3.4cm}@{}}
    \toprule
    \textbf{Option} & \textbf{Owner and first step} & \textbf{Evidence strength}\\
    \midrule
    % [Q5] PMEs value it, no detectable hiring penalty; [M16] about 1 firm in 5 offers it; [M17] offering goes
    % with a written FWA policy, an association; [M19] no firm group penalises it; [M20] design balance;
    % [S7] same hours over four longer days; [P21] TAFEP and SNEF resources; Lead decision 2: test, not proven
    \hyperlink{B-optA}{\textbf{\textcolor{Navy}{A\enspace Pilot:}}} put compressed weeks on the written menu and test them
      & SNEF and volunteer firms; NTUC: safeguards; TAFEP: guidance
      & Moderate; operating costs unmeasured\\ \addlinespace[2pt]
    % [M10, Q10] reservations where most roles need presence; [M12] q about .06 over 20 splits; EB s9 l.183:
    % preliminary association under a common pay-slope model; [P11] Job Redesign+ and the PTRG review for
    % flexi-load; [P5] flexi-load the weak spot; [R6, R7, G6] reduced load and part-time about pro-rata
    \hyperlink{B-optB}{\textbf{\textcolor{Navy}{B\enspace Support:}}} review which tasks need presence; flexi-load where fewer hours, not location, address the need
      & Firms, with Job Redesign+; MOM's PTRG review
      & Moderate; preliminary associations\\ \addlinespace[2pt]
    % [S7] 3-4 WFH days is an attribute level; [Q2, Q3] 31% and 23% at equal pay (appended; not printed);
    % [M13] no detectable penalty where neither presence need nor creativity concern; [V6, P7] trial abroad;
    % Lead decision 10: not a push for fully remote work
    \hyperlink{B-optC}{\textbf{\textcolor{Navy}{C\enspace Pilot:}}} 3--4-day WFH pilot in roles without presence needs % [S7] attribute level
      & Volunteer firms; independent evaluation
      & Exploratory\\ \addlinespace[2pt]
    % [P2] no administrative request counts; [P3] MOM survey: few in need ask formally; [B6] no why-not item;
    % [D7, D9] partial approvals, request item partly misread; draft items on B-optD-items
    \hyperlink{B-optD}{\textbf{\textcolor{Navy}{D\enspace Measure:}}} why-not-request and partial-approval survey items
      & MRSD checks the items against its survey, with NTUC and SNEF % fixer pass 7 Oct 2026 (mom_ps-01 b, mom_ps-02 b)
      & Gap documented\\
    \bottomrule
  \end{tabularx}\par}
  % Lead decision 10: no option needs a Guidelines change
  \Interp{None requires a change to the Guidelines.}
  \Backups{\BL{B-optA}{Option A} \BL{B-optD-items}{Survey items} \BL{B-pilot}{Pilot checklist}}
\end{frame}
\note{[Short path] So to the decision. We ask the partners to choose one measurement step and one bounded, voluntary practice pilot, and to nominate a lead for each. None requires a change to the Guidelines. % mandate: two-part decision; Lead decision 10
For the pilot, Option A is the most natural: SNEF could recruit volunteer firms to put compressed weeks on their written menu and test them, with NTUC co-designing safeguards, TAFEP providing guidance, and independent evaluation. The evidence is moderate, and operating costs are unmeasured. % [Q5, M16, M17, M19, M20]
Option B supports implementation: firms could review which tasks need presence, with Job Redesign+, and use flexi-load where fewer hours, not location, address the need. % [M10, Q10, P11, P5]
Option C, a three-to-four-day WFH pilot in roles without presence needs, is exploratory. % [M13, V6]
For measurement, Option D: MRSD, with NTUC and SNEF, could add why-not-request and partial-approval items where its surveys lack them; it adds no reporting duty for firms. % [P2, P3, B6, D9]
We close with three messages. % transition to M-conclusion
[Option letters link to B-optA to B-optD. Full deck: M-limits comes next.]}

% ==========================================================================
% 21. LIMITS   [full deck only]
% ==========================================================================
\begin{frame}[t,label=M-limits]{What we cannot yet say}
  % Six limits (ROUND5_SPEC s3 row 21; COMMENTS slide 21: add the pending free-slope check and the
  % 60 design cells). 10 pt table: limit | what it means.
  % Integrator, 7 Oct 2026: the third column (what would address it) is dropped; the spec asks for at most six
  % short limits, the next steps are on B-nextsteps and B-limits, and the slide was above the 80% density rule
  % (LEAD_DECISIONS_v3 rule 1).
  % Polish pass, 7 Oct 2026 (G8): the shared main-table style, \MainTableStyle = 10 pt, arraystretch 1.15, with
  % \RowSep between the limits; label column 3.0 cm (was 3.6), so the explanations need fewer lines.
  \FrameTop % first body token (centring; never a frame subtitle)
  {\MainTableStyle
  \begin{tabularx}{\linewidth}{@{}>{\raggedright\arraybackslash}p{3.0cm}>{\raggedright\arraybackslash}X@{}}
    \toprule
    \textbf{Limit} & \textbf{What it means}\\
    \midrule
    % [S5] hypothetical offers; [V1-V5] stated choices track directions better than levels; [S9] both questionnaires
    % said SNEF would use the results (fixer pass 7 Oct 2026, mom_ps-06 d)
    \textbf{Stated choices} & Hypothetical offers; respondents knew SNEF would use the results; rankings firmer than levels\\ \RowSep
    % [S4] unweighted online panels; [S3] only 12% of employer respondents' firms have 50 or fewer staff (not printed)
    % [S1, S2, S4] PMEs in or seeking PME roles and employers about PME hiring: no rank-and-file staff (worker-07)
    \textbf{Online panels} & Unweighted surveyed PMEs and firms, not national rates; no rank-and-file staff\\ \RowSep
    % [S2, S5] a hire between two equally qualified candidates; EB s9: hire, not request; operating costs not measured
    % [S2, S5; EB s9] a hiring reservation may also reflect what wanting the arrangement signals (editor-08)
    \textbf{Employer task} & Hiring, not a staff request; may also reflect candidate signals; operating costs unmeasured\\ \RowSep
    % respondent-clustered CIs likely too narrow with so few distinct situations (EB s9)
    % [EB s9] 10 versions x 6 tasks = 60 choice situations per survey; [S6, S6b] every pay gap 17-33%
    \textbf{Few design cells} & 60 choice situations per side; pay gaps 17--33\% by design; intervals likely too narrow\\ \RowSep
    % EB s9 l.183: no free-slope equal-pay check yet for the presence and creativity splits; [M12]
    \textbf{Employer splits} & Free-slope check pending: preliminary associations under a common pay-slope model\\ \RowSep
    % [D9] request item partly misread; [D7b] refusals only among reported requesters; [B6] no why-not item
    \textbf{Request item} & Partly misread; why PMEs do not ask is unmeasured\\
    \bottomrule
  \end{tabularx}\par}
  \Backups{\BL{B-limits}{All limits} \BL{B-nextsteps}{Next steps}}
\end{frame}
\note{[Full deck only] We acknowledge six limits. The choices are stated, so rankings are firmer than levels. % [S5, V1-V5]
Respondents knew SNEF would use the results, which may invite strategic answers on both sides. % [S9]
The panels are online and unweighted, so shares describe surveyed PMEs and firms, not national rates. % [S4]
The employer task is a hiring choice, not a reply to a staff request, and does not measure operating costs. % [S2, S5]
Each side saw only 60 choice situations, with pay gaps of 17 to 33 per cent, so intervals are likely too narrow. % [EB s9; S6, S6b]
The group-specific pay-slope check for the employer splits is pending, so they remain preliminary associations. % [EB s9 l.183, M12]
And the request item was partly misread. % [D9, B6]
MOM, SNEF and NTUC could each support a next step for these limits. With that, the decision we ask for. % implication (who could act) and transition to M-conclusion
}

% ==========================================================================
% 22. CONCLUSION: THE DECISION   [short path]
% ==========================================================================
\begin{frame}[t,label=M-conclusion]{The decision we ask of the partners}
  % Decision slide (mandate; COMMENTS slide 22): banner, then one take-away block with three robust
  % messages; the 490-PME detail is on B-followup; no "Thank you" here (M-thanks follows).
  % MANDATE: one bounded practice pilot; voluntary as on M-puzzle (fixer pass 7 Oct 2026: mom_ps-01 c, employer-04)
  \FrameTop
  \DecisionBanner{Decision sought today: one measurement step and one bounded, voluntary practice pilot, with a lead for each.}
  \vspace{0.20cm}% polish pass, 7 Oct 2026 (DESIGN_PLAN s2 W5): was 0.30 cm
  % Polish pass: no local font override in the block title; the restyled take-away card (preamble, G7) sets
  % the 11.5 pt bold Navy heading, the Navy top accent and the bold Navy numerals.
  \begin{alertblock}{Take-away messages}
    \begin{enumerate}\setlength{\itemsep}{0.12cm}
      % [E1, E2] 7.4 and 9.5 -> about 7-10% of pay, status I (indicative); [S6] gaps shown 17-33%; not a pay cut (mandate)
      \item PMEs strongly prefer hybrid work: the model implies \textbf{about 7--10\% of pay} (indicative, not a pay cut)
      % [Q5] 58% PMEs, 47% employers at equal pay (appended); [M19] no firm group penalises it (303-firm check); [M20] design balance; [S7] same hours, four longer days
      \item Compressed week: \textbf{most promising common-ground test}; evidence moderate; operating costs unmeasured
      % [R10] ranking robust: fully remote penalised most; [Q3] 23% at equal pay (appended; not printed); [Q16] 3-4 days vs remote p = .006 (303-firm); [S5] hiring choices
      \item Employers rank \textbf{fully remote last} in hiring: a ranking, not a salary price
    \end{enumerate}
  \end{alertblock}
  % mom_ps-01 (d): the proposed default, visible on the decision slide (Options D and A as on M-options)
  \Interp{We propose D (why-not and partial-approval items) and A (a compressed-week pilot).}
  \Backups{\BL{B-followup}{Follow-up} \BL{B-optD-items}{Survey items}}
\end{frame}
\note{[Short path] Three messages. First, PMEs strongly prefer hybrid work: the model implies about seven to ten per cent of pay, an indicative value, not a case for lower pay. % [E1, E2, S6]
Second, the compressed week, the same hours over four longer days, is the most promising common-ground test: PMEs value it and we detect no hiring penalty, but the evidence is moderate and operating costs are unmeasured. % [S7, Q5, M19, M20]
Third, employers rank fully remote work last when they hire; that is a ranking, not a salary price. % [R10, Q3, Q16]
We acknowledge that these are stated choices, and the employer splits remain preliminary. % [S4, S5; EB s9 l.183]
The policy problem is therefore fit, not aggregate provision: which arrangement, for which role, and what blocks its use. A measurement step would show why PMEs who need an arrangement do not ask; a bounded pilot would show whether the compressed week works in practice. % [B6, P2]
For example, MRSD could lead the measurement step, and SNEF a voluntary compressed-week pilot with volunteer firms, with NTUC co-designing safeguards. % owners as on M-options; worker-08: the since-clause on career concern [M23, B5] dropped to keep the short path under 18:00
We ask the partners today to choose one measurement step and one bounded, voluntary practice pilot, and to nominate a lead for each. % decision request (same sentence closes both versions)
[30-second close: PMEs strongly prefer hybrid work, at an indicative seven to ten per cent of pay. The compressed week is the most promising common-ground test, with moderate evidence and unmeasured operating costs, so any pilot needs worker safeguards. Employers rank fully remote last in hiring; that is a ranking, not a price. We ask the partners today to choose one measurement step and one bounded, voluntary practice pilot, and to nominate a lead for each.]
[Follow-up design and consent scope: B-followup.]}

